\corrigesChapitre{14}{Loi d’Ohm et puissance électrique}
\par\noindent\begin{minipage}{\linewidth}
\begin{corrige}{1}{Les symboles}
\begin{questions}
\item Tension et résistance.
\item Volt, V ; ohm, $\Omega$.
\end{questions}
\end{corrige}
\end{minipage}
\par\noindent\begin{minipage}{\linewidth}
\begin{corrige}{2}{Convertir}
\begin{questions}
\item $2\,200\unite\Omega$.
\item $0{,}015\unite A$.
\end{questions}
\end{corrige}
\end{minipage}
\par\noindent\begin{minipage}{\linewidth}
\begin{corrige}{3}{Trouver U}
\begin{questions}
\item $2{,}0\unite V$.
\item $U=RI$.
\end{questions}
\end{corrige}
\end{minipage}
\par\noindent\begin{minipage}{\linewidth}
\begin{corrige}{4}{Trouver I}
\begin{questions}
\item $I=0{,}020\unite A$.
\item $20\unite{mA}$.
\end{questions}
\end{corrige}
\end{minipage}
\par\noindent\begin{minipage}{\linewidth}
\begin{corrige}{5}{Trouver R}
\begin{questions}
\item $R=3/0{,}010=300\unite\Omega$.
\item Non : le quotient $0/0$ n’est pas défini.
\end{questions}
\end{corrige}
\end{minipage}
\par\noindent\begin{minipage}{\linewidth}
\begin{corrige}{6}{Une puissance}
\begin{questions}
\item $6\unite W$.
\item $6\unite J$ par seconde à puissance constante.
\end{questions}
\end{corrige}
\end{minipage}
\par\noindent\begin{minipage}{\linewidth}
\begin{corrige}{7}{Retrouver l’intensité}
\begin{questions}
\item $I=P/U=2\unite A$.
\item $1\unite A$.
\end{questions}
\end{corrige}
\end{minipage}
\par\noindent\begin{minipage}{\linewidth}
\begin{corrige}{8}{Une caractéristique}
\begin{questions}
\item $U/I=100\unite\Omega$ pour chaque point.
\item Vérifier l’alignement sur une droite proche de l’origine, avec dispersion compatible avec les mesures.
\end{questions}
\par\smallskip\noindent\textbf{Pour comprendre.}\enspace Former $1/0{,}01$, $2/0{,}02$ et $3/0{,}03$ : on obtient chaque fois 100 ohms. Le contrôle graphique porte aussi sur l’origine ; trois quotients cohérents ne valident pas des conditions jamais testées.\par
\pcFigureEx{figures/rev-cor-ohm-100.pdf}{Le quotient U/I vaut 100 ohms pour chacun des trois points ; le modèle est une droite passant par l’origine.}
\end{corrige}
\end{minipage}
\par\noindent\begin{minipage}{\linewidth}
\begin{corrige}{9}{Une limite dépassée}
\begin{questions}
\item $I=0{,}10\unite A$ ; $P=1{,}0\unite W$.
\item Non, la puissance calculée est quatre fois la limite ; ne pas réaliser ce montage.
\end{questions}
\par\smallskip\noindent\textbf{Pour comprendre.}\enspace $I=U/R=10/100=0{,}10\unite A$, puis $P=UI=1{,}0\unite W$. Comparer dans la même unité : $1{,}0/0{,}25=4$, donc la limite est dépassée.\par
\end{corrige}
\end{minipage}
\par\noindent\begin{minipage}{\linewidth}
\begin{corrige}{10}{À tension constante}
\begin{questions}
\item I est divisée par deux.
\item P est également divisée par deux dans ce modèle.
\end{questions}
\end{corrige}
\end{minipage}
\par\noindent\begin{minipage}{\linewidth}
\begin{corrige}{11}{À résistance constante}
\begin{questions}
\item I double.
\item P est multipliée par quatre, car U et I doublent.
\end{questions}
\end{corrige}
\end{minipage}
\par\noindent\begin{minipage}{\linewidth}
\begin{corrige}{12}{Une lampe réelle}
\begin{questions}
\item Non : les quotients diffèrent au-delà des éventuelles erreurs à examiner.
\item L’échauffement du filament peut modifier sa résistance.
\end{questions}
\pcFigureEx{figures/rev-cor-lampe-reelle.pdf}{Exemple de points illustrant les quotients 20 et 35 ohms donnés ; ce ne sont pas de nouvelles mesures de l’énoncé.}
\end{corrige}
\end{minipage}
\par\noindent\begin{minipage}{\linewidth}
\begin{corrige}{13}{Un graphique trompeur}
\begin{questions}
\item Non ; la pente I/U vaut $1/R$ pour le modèle ohmique.
\item Calculer $U/I$ avec I en A.
\end{questions}
\par\smallskip\noindent\textbf{Pour comprendre.}\enspace La pente se forme comme ordonnée/abscisse. Si I est vertical et U horizontal, elle vaut I/U en A/V. La résistance reste U/I en V/A, c’est-à-dire en ohms.\par
\pcFigureEx{figures/rev-cor-ohm-inverse.pdf}{Dans ces axes, la pente est I/U ; on retrouve R en prenant le quotient inverse U/I.}
\end{corrige}
\end{minipage}
\par\noindent\begin{minipage}{\linewidth}
\begin{corrige}{14}{Deux branches}
\begin{questions}
\item $I_1=0{,}060\unite A$ et $I_2=0{,}030\unite A$.
\item $P=6(0{,}060+0{,}030)=0{,}54\unite W$.
\end{questions}
\par\smallskip\noindent\textbf{Pour comprendre.}\enspace $I_1=6/100=0{,}060\unite A$, $I_2=6/200=0{,}030\unite A$. Le générateur fournit $I_G=0{,}090\unite A$, donc $P_G=UI_G=0{,}54\unite W$.\par
\pcFigureEx{figures/rev-cor-resistances-6.pdf}{Même tension de 6 V, courants de 60 mA et 30 mA ; la puissance totale vaut 6 × 0,090 = 0,54 W.}
\end{corrige}
\end{minipage}
\par\noindent\begin{minipage}{\linewidth}
\begin{corrige}{15}{Choisir un dipôle}
\begin{questions}
\item $R=250\unite\Omega$ ; $P=0{,}10\unite W$.
\item Oui, le calcul est inférieur à la limite ; la marge reste faible et les tolérances réelles doivent être vérifiées par le responsable du montage.
\end{questions}
\end{corrige}
\end{minipage}
\par\noindent\begin{minipage}{\linewidth}
\phantomsection\label{cor-defi-14}\pcSolutionDefi{14}{Une résistance pour le laboratoire}\begin{questions}
\item 100 Ω : 60 mA et 0,36 W. 220 Ω : environ 27,3 mA et 0,164 W.
\item 220 Ω respecte la limite ; 100 Ω la dépasse.
\item A en série, V aux bornes de la résistance, calibres et bornes vérifiés ; couper avant modification et rester dans les limites du matériel.
\end{questions}
\par\smallskip\noindent\textbf{Pour comprendre.}\enspace Pour chaque résistance, appliquer $I=U/R$ puis $P=UI$. Sous 6 V : $P_{100}=0{,}36\unite W>0{,}25\unite W$ ; $P_{220}\approx0{,}164\unite W<0{,}25\unite W$. Le contrôle de puissance précède la réalisation.\par
\pcFigureEx{figures/rev-cor-montage-ohm.pdf}{Mesurer U aux bornes de R et I en série ; le modèle retenu est ici 220 ohms sous 6 V.}
\end{minipage}
\par\noindent\begin{minipage}{\linewidth}
\phantomsection\label{cor-activite-14}\pcSousTitre{Pistes pour l’investigation}
Calculer U/I seulement pour I non nul, avec I en A. Un alignement proche de l’origine et des quotients sensiblement constants sont cohérents avec une résistance constante. Calculer P = UI au maximum de fonctionnement et comparer à la puissance admissible. Une dérive thermique peut modifier la résistance ; il faut respecter les limites avant de chercher une belle droite.\par\end{minipage}
